\documentclass[11pt,a4paper,oneside]{book}

% Encoding and fonts
\usepackage[utf8]{inputenc}
\usepackage[T1]{fontenc}
\usepackage{lmodern}

% Page geometry
\usepackage[
  top=2.5cm,
  bottom=2.5cm,
  left=2.5cm,
  right=2.5cm,
  headheight=14pt
]{geometry}

% Language
\usepackage[english]{babel}

% Colors
\usepackage{xcolor}
\definecolor{darkblue}{HTML}{2C3E50}
\definecolor{codebg}{HTML}{F5F5F5}
\definecolor{codeframe}{HTML}{DDDDDD}
\definecolor{linkcolor}{HTML}{2C3E50}
\definecolor{chaptercolor}{HTML}{2C3E50}

% Code listings
\usepackage{listings}
\lstset{
  backgroundcolor=\color{codebg},
  frame=single,
  rulecolor=\color{codeframe},
  basicstyle=\ttfamily\small,
  breaklines=true,
  breakatwhitespace=false,
  tabsize=4,
  showstringspaces=false,
  numbers=none,
  xleftmargin=0.5em,
  xrightmargin=0.5em,
  aboveskip=1em,
  belowskip=1em,
  columns=fullflexible,
  keepspaces=true,
  literate={á}{{\'a}}1 {é}{{\'e}}1 {í}{{\'i}}1 {ó}{{\'o}}1 {ú}{{\'u}}1
}

% Headers and footers
\usepackage{fancyhdr}
\pagestyle{fancy}
\fancyhf{}
\fancyhead[L]{\small\itshape\leftmark}
\fancyhead[R]{\small\thepage}
\fancyfoot{}
\renewcommand{\headrulewidth}{0.4pt}

% Plain pages (chapter starts)
\fancypagestyle{plain}{
  \fancyhf{}
  \fancyfoot[C]{\small\thepage}
  \renewcommand{\headrulewidth}{0pt}
}

% Chapter and section styling
\usepackage{titlesec}

\titleformat{\chapter}[display]
  {\normalfont\huge\bfseries\color{chaptercolor}}
  {\chaptertitlename\ \thechapter}{20pt}{\Huge}
\titlespacing*{\chapter}{0pt}{-20pt}{30pt}

\titleformat{\section}
  {\normalfont\Large\bfseries\color{chaptercolor}}
  {\thesection}{1em}{}

\titleformat{\subsection}
  {\normalfont\large\bfseries\color{chaptercolor}}
  {\thesubsection}{1em}{}

% Links
\usepackage[
  colorlinks=true,
  linkcolor=linkcolor,
  urlcolor=linkcolor,
  citecolor=linkcolor,
  bookmarks=true,
  bookmarksnumbered=true,
  pdfstartview=FitH
]{hyperref}

% Tables
\usepackage{longtable}
\usepackage{booktabs}
\usepackage{array}

% Graphics
\usepackage{graphicx}

% Pandoc code highlighting compatibility
\usepackage{framed}
\usepackage{fancyvrb}
\newenvironment{Shaded}{\begin{snugshade}}{\end{snugshade}}
\definecolor{shadecolor}{HTML}{F5F5F5}
\DefineVerbatimEnvironment{Highlighting}{Verbatim}{
  commandchars=\\\{\},
  fontsize=\small,
  xleftmargin=0.5em
}
\newcommand{\KeywordTok}[1]{\textcolor[HTML]{007020}{\textbf{#1}}}
\newcommand{\DataTypeTok}[1]{\textcolor[HTML]{902000}{#1}}
\newcommand{\DecValTok}[1]{\textcolor[HTML]{40A070}{#1}}
\newcommand{\BaseNTok}[1]{\textcolor[HTML]{40A070}{#1}}
\newcommand{\FloatTok}[1]{\textcolor[HTML]{40A070}{#1}}
\newcommand{\ConstantTok}[1]{\textcolor[HTML]{880000}{#1}}
\newcommand{\CharTok}[1]{\textcolor[HTML]{4070A0}{#1}}
\newcommand{\SpecialCharTok}[1]{\textcolor[HTML]{4070A0}{#1}}
\newcommand{\StringTok}[1]{\textcolor[HTML]{4070A0}{#1}}
\newcommand{\VerbatimStringTok}[1]{\textcolor[HTML]{4070A0}{#1}}
\newcommand{\SpecialStringTok}[1]{\textcolor[HTML]{BB6688}{#1}}
\newcommand{\ImportTok}[1]{#1}
\newcommand{\CommentTok}[1]{\textcolor[HTML]{60A0B0}{\textit{#1}}}
\newcommand{\DocumentationTok}[1]{\textcolor[HTML]{BA2121}{\textit{#1}}}
\newcommand{\AnnotationTok}[1]{\textcolor[HTML]{60A0B0}{\textbf{\textit{#1}}}}
\newcommand{\CommentVarTok}[1]{\textcolor[HTML]{60A0B0}{\textbf{\textit{#1}}}}
\newcommand{\OtherTok}[1]{\textcolor[HTML]{007020}{#1}}
\newcommand{\FunctionTok}[1]{\textcolor[HTML]{06287E}{#1}}
\newcommand{\VariableTok}[1]{\textcolor[HTML]{19177C}{#1}}
\newcommand{\ControlFlowTok}[1]{\textcolor[HTML]{007020}{\textbf{#1}}}
\newcommand{\OperatorTok}[1]{\textcolor[HTML]{666666}{#1}}
\newcommand{\BuiltInTok}[1]{#1}
\newcommand{\ExtensionTok}[1]{#1}
\newcommand{\PreprocessorTok}[1]{\textcolor[HTML]{BC7A00}{#1}}
\newcommand{\AttributeTok}[1]{\textcolor[HTML]{7D9029}{#1}}
\newcommand{\RegionMarkerTok}[1]{#1}
\newcommand{\InformationTok}[1]{\textcolor[HTML]{60A0B0}{\textbf{\textit{#1}}}}
\newcommand{\WarningTok}[1]{\textcolor[HTML]{60A0B0}{\textbf{\textit{#1}}}}
\newcommand{\AlertTok}[1]{\textcolor[HTML]{FF0000}{\textbf{#1}}}
\newcommand{\ErrorTok}[1]{\textcolor[HTML]{FF0000}{\textbf{#1}}}
\newcommand{\NormalTok}[1]{#1}

% Tight list (pandoc compatibility)
\providecommand{\tightlist}{%
  \setlength{\itemsep}{0pt}\setlength{\parskip}{0pt}}

% Copyleft symbol (mirrored ©)
\providecommand{\textcopyleft}{\reflectbox{\copyright}}


% No paragraph indent, space between paragraphs
\setlength{\parindent}{0pt}
\setlength{\parskip}{0.6em}

% Quote environment styling (for epigraph)
\usepackage{etoolbox}
\AtBeginEnvironment{quote}{%
  \itshape\color{darkblue}%
}

% PDF metadata
\hypersetup{
  pdftitle={Ten Brief Lessons on Web Programming},
  pdfauthor={Giuseppe Costanzi},
  pdfsubject={Web Programming},
  pdfkeywords={PHP, JavaScript, HTML, CSS, SQLite, web development}
}

\begin{document}

% --- COVER PAGE ---
\begin{titlepage}
\newgeometry{margin=0pt}
\pagecolor{darkblue}
\begin{center}

\vspace*{6cm}

{\fontsize{28}{34}\selectfont\bfseries\color{white}
Ten Brief Lessons\\[0.4cm]on Web Programming}

\vspace{0.8cm}

{\large\color{white!80}
A concept notebook on web programming}

\vspace{2cm}

{\small\color{white!60}\itshape
``There are 10 types of people in the world:\\
those who understand binary, and those who don't.''}

\vspace{0.3cm}

{\small\color{white!80}
Which side are you on?}

\vspace{\fill}

{\large\color{white}
Giuseppe Costanzi}

\vspace{2cm}

\end{center}
\restoregeometry
\nopagecolor
\end{titlepage}

% --- BLANK PAGE ---
\thispagestyle{empty}
\vspace*{\fill}
\begin{center}
\textit{\color{gray}This page is intentionally left blank.}
\end{center}
\vspace*{\fill}
\newpage

% --- COPYRIGHT / COLOPHON ---
\thispagestyle{empty}
\vspace*{\fill}

{\small

\textbf{Ten Brief Lessons on Web Programming}\\
\textit{A concept notebook on web programming}

\vspace{1em}

Copyleft \textcopyleft\ 2026 Giuseppe Costanzi \textit{aka} \texttt{1966bc}

\vspace{1em}

Version 1.0 --- First edition, MMXXVI AD / MMDCCLXXIX \textit{ab urbe condita}

\vspace{1em}

This book is free software: you can redistribute it and/or modify
it under the terms of the GNU General Public License as published by
the Free Software Foundation, either version 3 of the License, or
(at your option) any later version.

\vspace{1em}

Source code: \url{https://github.com/1966bc/bibliotheca}

\vspace{1em}

Built with pure PHP, JavaScript, HTML, CSS and SQLite.\\
Typeset with \LaTeX\ via Pandoc.\\
No frameworks were harmed in the making of this book.

\vspace{1em}

\texttt{giuseppe.costanzi@gmail.com}

\vspace{1em}

{\footnotesize\itshape
Milky Way Galaxy, Solar System, third planet from the Sun\\
41\textdegree{}53'35''N \enspace 12\textdegree{}28'58''E --- Roma, Italia}
}

\newpage

% --- DEDICATION ---
\thispagestyle{empty}
\vspace*{5cm}
\begin{center}
\itshape\large
To the next generation of programmers:\\
learn to read, write and compute ---\\
before it's too late.
\end{center}
\newpage

% --- TABLE OF CONTENTS ---
{
% Add some stretch to help LaTeX distribute TOC entries better
\setlength{\parskip}{0.1em plus 0.1em}
\tableofcontents
}
\newpage

% --- BODY ---
% Start chapter numbering from 0
\setcounter{chapter}{-1}

$body$

\end{document}
